% load packages
%\usepackage{multicol}
\usepackage{fontspec}
\usepackage{xltxtra}
%\usepackage{xcolor}
\usepackage{listings}
\usepackage{fvextra}
% babel: Beim gewoehnlichen "pdf:"-Format haengt Quarto bei aktivem
% "lang: de" automatisch \usepackage[...]{babel} mit der richtigen
% Sprachoption ins Template ein. Die titlepage-pdf-Extension (aktives
% Format, s. _quarto.yml) ersetzt aber das Template-Partial pandoc.tex
% komplett und uebernimmt diesen $if(lang)$-Block NICHT mit -- ohne babel
% hat csquotes (s. u.) keine Sprachinfo und faellt auf den englischen
% Anfuehrungszeichen-Stil ("...") zurueck statt „..." zu verwenden. Deshalb
% hier manuell geladen. Bei einem Formatwechsel weg von titlepage-pdf
% pruefen, ob das zu einer doppelten babel-Ladung fuehrt (dann diese Zeile
% wieder auskommentieren).
\usepackage[ngerman]{babel}
% csquotes wird zusammen mit der Metadaten-Variable "csquotes: true" im
% titlepage-pdf-Formatblock von _quarto.yml gebraucht: Pandoc gibt
% Anfuehrungszeichen aus geraden "..." dann als \enquote{...} statt als
% literale Unicode-Zeichen aus, und csquotes waehlt anhand der von babel
% geladenen Sprache automatisch den passenden deutschen Stil („...") statt
% des sonst verwendeten englischen ("..."). Das Paket alleine (ohne die
% "csquotes: true"-Variable) haette keinen Effekt -- Pandoc muesste sonst
% weiterhin literale Zeichen statt \enquote{} ausgeben.
\usepackage{csquotes}
\usepackage{tocloft}  % more space in TOC
\usepackage{microtype}
% soul stellt \st{...} (Durchstreichen) bereit, das Pandoc fuer Markdown
% ~~...~~ (Strikeout) im LaTeX-Output verwendet. Der normale Pandoc/Quarto-
% Standardtemplate (latex.common) laedt soul selbst automatisch, wenn ein
% Dokument Strikeout enthaelt -- die titlepage-pdf-Extension (s. Format-
% Block "titlepage-pdf" in _quarto.yml) ersetzt aber template-partials wie
% pandoc.tex komplett und uebernimmt diese Bedingung nicht mit, wodurch
% \st{...} sonst als "Undefined control sequence" fehlschlaegt. Deshalb
% hier unconditional geladen (unabhaengig davon, ob ~~...~~ vorkommt).
\usepackage{soul}

% Farbige Emoji-Darstellung im PDF: Der Buchtext enthaelt literale Unicode-
% Emoji, auch in Ueberschriften (-> PDF-Bookmarks/hyperref) und in
% R-Code-Kommentaren (-> Highlighting/Verbatim-Umgebung). mainfont/sansfont/
% monofont enthalten keine Emoji-Glyphen. Registriert hier nur eine
% luaotfload-Fallback-Schriftklasse ("emojifallback"); tatsaechlich
% aktiviert wird sie ueber die RawFeature={fallback=emojifallback}-Option
% an mainfontoptions/sansfontoptions/monofontoptions (siehe _quarto.yml,
% pdf/titlepage-pdf-Format). Das ist bewusst KEIN Text-/Zeichenketten-Eingriff
% (z. B. per LuaTeX process_input_buffer-Callback), weil ein solcher Eingriff
% Zeichenfolgen roh im Eingabepuffer umschreibt und dabei den TOC/Bookmark-
% Text (hyperref/pdfstringdef) sowie den catcode-sensitiven Inhalt von
% Highlighting/Verbatim-Bloecken beschaedigen kann -- getestet und
% verworfen. Der reine Glyph-Fallback wirkt dagegen unabhaengig vom Kontext
% korrekt ueberall dort, wo die aktuelle Schrift die Glyphe nicht hat.
% Bekannter, akzeptierter Kompromiss: mehrteilige ZWJ-Emoji (z. B. Person +
% Beruf + Geschlecht wie "🧑‍🏫") werden dadurch als mehrere nebeneinander
% stehende Einzelsymbole gesetzt statt als ein fusioniertes Emoji-Glyph --
% inhaltlich weiterhin eindeutig erkennbar, nur nicht perfekt verschmolzen.
% Zweite Fallback-Stufe "Noto Sans Symbols2": deckt Zeichen ab, die Noto
% Color Emoji nicht enthaelt, aber trotzdem im Buchtext vorkommen -- z. B.
% Wuerfelaugen U+2680-2685 und Check Mark/Ballot X U+2713/U+2717. Das sind
% keine Farb-Emoji, sondern gewoehnliche Symbol-Glyphen, die in einer
% reinen Symbol-Schrift stehen; check_emoji_coverage.py (siehe Skill
% textlayout-ueberarbeitung) hat das per lualatex-Testkompilat verifiziert.
\directlua{
  luaotfload.add_fallback
    ("emojifallback",
     { "Noto Color Emoji:mode=harf;",
       "Noto Sans Symbols2:mode=harf;" }
    )
}


\definecolor{ycol}{RGB}{230,159,0}
\definecolor{modelcol}{RGB}{86,180,233}
\definecolor{errorcol}{RGB}{0,158,115}
\definecolor{beta0col}{RGB}{213,94,0}
\definecolor{beta1col}{RGB}{0,114,178}
\definecolor{xcol}{RGB}{204,121,167}


% Verbesserte Silbentrennung
\lefthyphenmin=2
\righthyphenmin=3
\hyphenpenalty=50
\exhyphenpenalty=50

% Für sehr lange Wörter
\tolerance=1000
\emergencystretch=3em
\hbadness=10000


\lstset{
  breaklines=true
}

\DefineVerbatimEnvironment{Highlighting}{Verbatim}{breaklines,commandchars=\\\{\}}
\DefineVerbatimEnvironment{OutputCode}{Verbatim}{breaklines,commandchars=\\\{\}}


\raggedbottom
\flushbottom


\setlength{\cftchapnumwidth}{4em} % Increase this to give more room to the number



% rm page number of "part pages", added 2025-04-17
\makeatletter
\def\@part[#1]#2{%
  \ifnum \c@secnumdepth > -1
    \refstepcounter{part}%
    \addcontentsline{toc}{part}{\thepart\hspace{1em}#1}%
  \else
    \addcontentsline{toc}{part}{#1}%
  \fi
  \markboth{}{}%
  \clearpage
  \thispagestyle{empty} % <-- removes page number
  {\centering
    \interlinepenalty \@M
    \normalfont
    \ifnum \c@secnumdepth > -1
      \huge\bfseries \partname~\thepart
      \par\nobreak
      \vskip 20pt
    \fi
    \Huge \bfseries #2\par
    \vskip 40pt
  }
  \@afterheading
}
\makeatother


